% !TeX root = ../regression_logistique.tex
\chapter{Notes brutes et standardisation}
\label{ch:standardisation}

\orient{%
justifier le choix d'une valeur d'imputation ; calculer une cote standard ;
délimiter ce que la standardisation conserve du nuage ; dire sur quelles lignes
ces statistiques doivent être ajustées}{%
moyenne, médiane, écart-type (\annexeref{ann:prerequis}) ;
repère~\ref{ch:probleme}}{%
essentiel}

\section{Fichier d'entraînement}

\texttt{dataset\_train.csv} contient $1\,600$ lignes, une par élève. Les six
premières colonnes donnent l'index, la maison, le nom, la date de naissance et la
main directrice ; les treize suivantes donnent une note par matière. La maison
est renseignée pour les $1\,600$ élèves : c'est la réponse que le modèle doit
retrouver.

Deux d'entre elles servent ici, \Herb et \Runes : leurs distributions par maison
se recouvrent peu, et elles suffisent à isoler les \Gryf (ch.~\ref{ch:nuage}).
Le modèle complet en retient dix. Ce choix de deux matières a été fait après
examen du fichier étiqueté entier ; sa portée méthodologique est discutée au
repère~\ref{ch:probleme}.

\begin{table}[htbp]
\centering
\begin{tabular}{lrrrrrr}
\toprule
matière & absentes & minimum & maximum & médiane & moyenne & écart-type \\
\midrule
\Herb     & $33$ & $-10{,}30$ & $11{,}61$  & $3{,}469$   & $1{,}189$   & $5{,}175$ \\
\Runes & $35$ & $283{,}87$ & $745{,}40$ & $463{,}918$ & $495{,}052$ & $105{,}186$ \\
\bottomrule
\end{tabular}
\caption{Les deux colonnes employées, telles qu'elles figurent dans le fichier.
La moyenne et l'écart-type sont calculés après remplacement des notes absentes,
donc sur les $1\,600$ lignes.}
\label{tab:brut}
\end{table}

\section{Notes absentes}

Certaines cases du fichier sont vides. Le score étant une somme de produits, il
exige une valeur pour chaque terme : une case vide arrête le traitement de la
ligne entière.

Deux traitements existent. Écarter les lignes incomplètes coûte de l'effectif et
suppose que les lignes écartées ressemblent aux lignes conservées, hypothèse
invérifiable. Attribuer une valeur à chaque case vide --- l'\emph{imputation} ---
fabrique des mesures. Ce cours retient l'imputation par la médiane : coupant la
colonne en deux moitiés de même effectif, elle résiste au déplacement des valeurs
extrêmes, que la moyenne suit.

Toute imputation par une constante concentre les lignes incomplètes en un point
du nuage et resserre la dispersion de la colonne. Sa validité tient à deux
conditions : des cases vides rares, et une absence indépendante de la valeur
qu'elles auraient prise. La seconde reste invérifiable sur les seules
données~\cite{littlerubin2019}.

\begin{chiffres}[Cases vides des deux colonnes]
\begin{center}
\begin{tabular}{lrrl}
\toprule
 & \Herb & \Runes & \\
\midrule
notes absentes    & $33$ & $35$ & soit $68$ élèves, $4{,}25\,\%$ du fichier \\
médiane attribuée & $3{,}469$ & $463{,}918$ & \\
\addlinespace
moyenne avant     & $1{,}141$ & $495{,}748$ & sur les notes présentes \\
moyenne après     & $1{,}189$ & $495{,}052$ & elle monte d'un côté, baisse de l'autre \\
\bottomrule
\end{tabular}
\end{center}
Le remplissage déplace chaque moyenne de moins d'un centième d'écart-type :
$0{,}009$ pour \Herb, $0{,}007$ pour \Runes. Les $68$ élèves concernés le sont
chacun sur une seule des deux colonnes.
\end{chiffres}

\begin{releve}
L'élève $21$, dont la note d'\Runes est absente, reçoit $463{,}918$, soit
$-0{,}296$ une fois standardisée.
\end{releve}

\section{Différence d'échelle entre les colonnes}

Les deux colonnes mesurent une performance scolaire, mais dans des unités sans
rapport entre elles. Le score les additionne pourtant, après avoir multiplié
chacune par son coefficient. À coefficient égal, une matière dont les notes
s'étalent sur une large plage pèse alors davantage qu'une autre, pour la seule
raison de son échelle de notation. Deux conséquences en découlent : les coefficients
cessent d'être comparables entre eux, et la distance entre deux élèves du plan
dépend du barème de chaque matière plutôt que de leurs résultats.

\begin{releve}
Une note d'\Herb se lit entre $-10$ et $+12$, une note d'\Runes autour de
$500$, et les secondes sont vingt fois plus dispersées : dans $w_1x_1+w_2x_2$, un
$w_2$ de $1$ produit donc un effet vingt fois plus grand qu'un $w_1$ de $1$.
\end{releve}

Ramener les deux colonnes à une échelle commune demande de fixer une origine et
une unité pour chacune d'elles.

Le \emph{centrage} soustrait la moyenne $\mu$ : l'écart $v-\mu$ situe chaque
note par rapport à la promotion, avec des valeurs positives au-dessus de la
moyenne et négatives en dessous.

La \emph{réduction} divise cet écart par l'écart-type $\sigma$, mesure de la
dispersion des notes autour de leur moyenne. Chaque valeur s'exprime alors en
nombre d'écarts-types.

\[
  x=\frac{\annup{accent}{v-\mu}{\annl{écart à la moyenne}\\ \annl{de la colonne}}}
         {\ann{warning}{\sigma}{\annl{écart-type}\\ \annl{de la colonne}}}
\]

Le résultat porte le nom de \emph{cote standard}, ou \emph{$z$-score} en
anglais ; on le note ici $x$, la lettre $z$ étant réservée au score du modèle
(ch.~\ref{ch:nuage}). Par construction, la colonne transformée a pour moyenne $0$
et pour écart-type $1$, quelle que soit la matière d'origine. Une valeur
standardisée se lit alors directement : $+1$ désigne un élève situé une
dispersion typique au-dessus de sa moyenne, en \Herb comme en \Runes, et les deux
colonnes deviennent comparables terme à terme.

\begin{figure}[htbp]
\centering
\begin{graphique}{Deux règles graduées superposées, reliées par des flèches verticales. En haut, la note brute d'Herbology de -10,30 à 11,61, moyenne à 1,189. En bas, la même note standardisée, de -2,22 à +2,01, zéro à la position de la moyenne. L'élève 3, en rouge, passe de -6,497 à -1,485.}
\begin{tikzpicture}[font=\small,
  grad/.style={ink!45},
  lab/.style={font=\scriptsize,text=ink},
  ruler/.style={font=\footnotesize,text=ink,anchor=west}]
% Position sur la regle : (v + 11) / 23 * 12, pour v de -11 a +12.
\node[ruler] at (0,1.05) {\Herb, note brute};
\draw[ink!45] (0,0) -- (12,0);
\foreach \v/\lab in {0.365/{$-10{,}30$},6.359/{$1{,}189$},11.797/{$11{,}61$}}{
  \draw[grad] (\v,-0.12) -- (\v,0.12);
  \node[lab,anchor=south] at (\v,0.18) {\lab};
}
\draw[warning] (2.349,-0.12) -- (2.349,0.12);
\node[lab,text=warning,anchor=south] at (2.349,0.18) {$-6{,}497$};

\draw[-{Latex[length=2mm]},ink!45] (6.359,-0.28) -- (6.359,-1.72);
\draw[-{Latex[length=2mm]},warning!70] (2.349,-0.28) -- (2.349,-1.72);

\node[ruler,text=accent] at (0,-2.95) {\Herb, note standardisée};
\draw[accent!60] (0,-2) -- (12,-2);
\foreach \v/\lab in {0.365/{$-2{,}22$},6.359/{$0$},11.797/{$+2{,}01$}}{
  \draw[accent!60] (\v,-2.12) -- (\v,-1.88);
  \node[lab,text=accent,anchor=north] at (\v,-2.18) {\lab};
}
\draw[warning] (2.349,-2.12) -- (2.349,-1.88);
\node[lab,text=warning,anchor=north] at (2.349,-2.18) {$-1{,}485$};
\end{tikzpicture}
\end{graphique}
\caption{Sur un seul axe, la transformation est affine croissante : elle déplace
l'origine et change l'unité, sans modifier l'ordre des élèves. En rouge,
l'élève~$3$.}
\label{fig:graduation}
\end{figure}
\FloatBarrier

Chaque axe est divisé par un nombre différent --- $5{,}175$ d'un côté,
$105{,}186$ de l'autre. La transformation conserve donc la structure affine du
nuage : l'ordre sur chaque coordonnée, l'alignement des points, la position
intermédiaire d'un point entre deux autres, et la correspondance terme à terme
entre un modèle affine sur les notes standardisées et un modèle affine sur les
notes brutes.

Distances euclidiennes, rapports de distances et angles relèvent en revanche du
plan standardisé seul, où les deux axes portent la même unité. Les notions
d'angle, de perpendiculaire et de distance employées au chapitre~\ref{ch:nuage}
s'entendent dans ce plan. La standardisation installe une géométrie, choisie
parce qu'elle rend les coefficients comparables entre eux.

\begin{figure}[htbp]
\centering
\begin{graphique}{Deux nuages de points côte à côte, mêmes 534 élèves. À gauche, notes brutes : l'axe horizontal va de -12 à 13, l'axe vertical de 270 à 760. À droite, notes standardisées : les deux axes vont de -2,8 à 2,8. Dans les deux, les Gryffindor en disques bleus occupent la zone en haut à gauche, les autres maisons en croix orangées le reste.}
\begin{tikzpicture}[baseline=(current bounding box.north),
  alt={Nuage des notes brutes de 534 élèves. Herbology va de -12 à 13 et Ancient Runes de 270 à 760 ; les Gryffindor se concentrent en haut à gauche.}]
\begin{axis}[width=.42\textwidth,height=6.2cm,
  axis lines=box,axis line style={ink!35},
  grid=major,grid style={gridgray!35},
  tick label style={font=\scriptsize},label style={font=\scriptsize},
  title style={font=\small},title={notes brutes},
  xlabel={\Herb},ylabel={\Runes},
  xmin=-12,xmax=13,ymin=270,ymax=760,ytick={300,450,600,750}]
  \addplot[only marks,mark=x,mark size=1.3pt,line width=.35pt,warning!85,opacity=.55]
    table[x=x,y=y] {data/brut_autres.dat};
  \addplot[only marks,mark=*,mark size=1.1pt,accent,opacity=.75]
    table[x=x,y=y] {data/brut_gryffindor.dat};
\end{axis}
\end{tikzpicture}\hfill
\begin{tikzpicture}[baseline=(current bounding box.north),
  alt={Même nuage après standardisation. Les deux axes vont de -2,8 à 2,8 ; la position relative des points et la séparation des groupes sont conservées.}]
\begin{axis}[width=.42\textwidth,height=6.2cm,
  axis lines=box,axis line style={ink!35},
  grid=major,grid style={gridgray!35},
  tick label style={font=\scriptsize},label style={font=\scriptsize},
  title style={font=\small},title={notes standardisées},
  xlabel={\Herb, en écarts-types},
  ylabel={\Runes, en écarts-types},
  xmin=-2.8,xmax=2.8,ymin=-2.8,ymax=2.8,xtick={-2,0,2},ytick={-2,0,2}]
  \addplot[only marks,mark=x,mark size=1.3pt,line width=.35pt,warning!85,opacity=.55]
    table[x=x,y=y] {data/autres.dat};
  \addplot[only marks,mark=*,mark size=1.1pt,accent,opacity=.75]
    table[x=x,y=y] {data/gryffindor.dat};
\end{axis}
\end{tikzpicture}
\end{graphique}
\caption{$534$ élèves : \Gryf en bleu, autres maisons en orangé. Les
transformations affines des axes préservent l'ordre et les alignements ; les
distances et les angles changent.}
\label{fig:brut-vs-standard}
\end{figure}
\FloatBarrier

\section{Application à un élève}

L'élève $3$ a $-6{,}497$ en \Herb et $523{,}982$ en \Runes.

\[
  x_1=\frac{-6{,}497-1{,}189}{5{,}175}=-1{,}485,
  \qquad
  x_2=\frac{523{,}982-495{,}052}{105{,}186}=+0{,}275 .
\]

Ces deux valeurs sont les coordonnées du point qui représente cet élève dans le
plan $(x_1,x_2)$. La première le situe à un écart-type et demi au-dessous de la
moyenne en \Herb, la seconde à un quart d'écart-type au-dessus en \Runes.

\begin{resultat}[Statistiques à conserver]
La médiane, la moyenne et l'écart-type sont ajustés une fois, sur les seules
lignes d'apprentissage, et font partie du modèle au même titre que les
coefficients. Une observation nouvelle se standardise avec ces trois nombres
figés (ch.~\ref{ch:predire}).
\end{resultat}
